% ==============================================================================
% Fichier     : style.sty
% Titre       : Feuille de style académique générique
% Auteur      : MAGUENA NDENE Marcel Allan
% Institution : Université de Yaoundé I — ENSP
% Description : Feuille de style réutilisable pour tout devoir/rapport.
%               Toutes les métadonnées sont paramétrables depuis main.tex.
%               Palette : bleu académique / orange accent / gris neutres.
% Usage       : \usepackage{style} dans main.tex, puis définir les
%               commandes de métadonnées listées en section I.
% Révision    : 2026
% ==============================================================================

\ProvidesPackage{style}[2026/03/01 Feuille de style académique — MAGUENA NDENE Marcel Allan]

% ==============================================================================
% I. PACKAGES REQUIS
% ==============================================================================

\RequirePackage[T1]{fontenc}
\RequirePackage[utf8]{inputenc}
\RequirePackage[provide=*,french]{babel}
\RequirePackage{lmodern}
\RequirePackage[protrusion=true,expansion=false]{microtype}

% --- Mise en page ---
\RequirePackage[
    top=2.5cm, bottom=2.5cm,
    left=2.8cm, right=2.5cm,
    headheight=14pt
]{geometry}

% --- Couleurs ---
\RequirePackage[dvipsnames, table]{xcolor}

% --- Graphiques & Figures ---
\RequirePackage{graphicx}
\RequirePackage{float}
\RequirePackage{wrapfig}
\RequirePackage{tikz}
\usetikzlibrary{shapes.geometric, arrows.meta, positioning, shadows, calc}

% --- Tableaux ---
\RequirePackage{booktabs}
\RequirePackage{longtable}
\RequirePackage{tabularx}
\RequirePackage{multirow}
\RequirePackage{array}
\RequirePackage{colortbl}

% --- Mathématiques ---
\RequirePackage{amsmath}
\RequirePackage{amssymb}
\RequirePackage{amsthm}

% --- Environnements décoratifs ---
\RequirePackage[most]{tcolorbox}
\tcbuselibrary{breakable, skins, theorems}

% --- Listes ---
\RequirePackage{enumitem}

% --- Code source ---
\RequirePackage{listings}
\RequirePackage{lstautogobble}

% --- En-têtes et pieds de page ---
\RequirePackage{fancyhdr}

% --- Hyperréférences (couleurs héritent de la palette) ---
\RequirePackage[
    colorlinks=true,
    linkcolor=BleuPrimaire,
    urlcolor=BleuSecondaire,
    citecolor=OrangeAccent,
    bookmarks=true,
    bookmarksnumbered=true
]{hyperref}

% --- Bibliographie ---
\RequirePackage[
    backend=biber,
    style=ieee,
    sorting=nyt
]{biblatex}

% --- Index ---
\RequirePackage{makeidx}
\makeindex

% --- Glossaires ---
\RequirePackage[toc, acronym]{glossaries}

% --- Divers ---
\RequirePackage{epigraph}
\RequirePackage{pifont}
\RequirePackage{marvosym}
\RequirePackage{setspace}
\RequirePackage{pdflscape}
\RequirePackage{afterpage}
\RequirePackage{changepage}


% ==============================================================================
% II. MÉTADONNÉES PARAMÉTRABLES
% ==============================================================================
%
% Définissez ces commandes AVANT \begin{document} dans votre main.tex.
% Des valeurs par défaut sont fournies ici pour éviter les erreurs si
% une commande est oubliée.
%
% Exemple dans main.tex :
%   \renewcommand{\doctitre}{Analyse de Risque et Gestion de Crise}
%   \renewcommand{\docsoustitre}{Gouvernance, Conformité, Résilience}
%   \renewcommand{\doctypedevoir}{Rapport de TP}
%   \renewcommand{\docue}{GES4062}
%   \renewcommand{\doclogo}{images/logo.png}
%   \renewcommand{\docauteur}{MAGUENA NDENE Marcel Allan}
%   \renewcommand{\docinstitution}{Université de Yaoundé I — ENSP}
%   \renewcommand{\docdepartement}{Département de Génie Informatique}
%   \renewcommand{\docanneescolaire}{2025--2026}
%   \renewcommand{\docnormes}{ISO 27001 · ISO 19011 · ISO 22301}
%   \renewcommand{\docepigraphe}{Votre citation ici.}
% ==============================================================================

\newcommand{\doctitre}{Titre du devoir}
\newcommand{\docsoustitre}{Sous-titre ou thématique}
\newcommand{\doctypedevoir}{Rapport}
\newcommand{\docue}{UE-XXXX}
\newcommand{\doclogo}{logo.png}
\newcommand{\docauteur}{Auteur}
\newcommand{\docinstitution}{Institution}
\newcommand{\docdepartement}{Département}
\newcommand{\docanneescolaire}{20XX--20XX}
\newcommand{\docnormes}{}           % Laisser vide si non applicable
\newcommand{\docepigraphe}{}        % Laisser vide pour ne pas afficher d'épigraphe


% ==============================================================================
% III. PALETTE CHROMATIQUE — BLEU ACADÉMIQUE + ORANGE + GRIS
% ==============================================================================

% ── Bleus (hiérarchie visuelle principale) ───────────────────────────────────
\definecolor{BleuPrimaire}{RGB}{0, 56, 120}          % Bleu marine — chapitres, titres
\definecolor{BleuSecondaire}{RGB}{0, 102, 179}       % Bleu vif — sections, liens
\definecolor{BleuClair}{RGB}{210, 228, 250}          % Bleu très pâle — fonds de boîtes
\definecolor{BleuMoyen}{RGB}{100, 160, 220}          % Bleu intermédiaire — filets légers

% ── Oranges d'accent ─────────────────────────────────────────────────────────
\definecolor{OrangeAccent}{RGB}{210, 100, 0}         % Orange foncé — accent fort
\definecolor{OrangePale}{RGB}{255, 237, 210}         % Orange très pâle — fond alerte
\definecolor{OrangeMoyen}{RGB}{240, 160, 50}         % Orange moyen

% ── Gris neutres ─────────────────────────────────────────────────────────────
\definecolor{NoirProfond}{RGB}{15, 15, 15}           % Quasi-noir — en-têtes
\definecolor{GrisAncre}{RGB}{45, 45, 45}             % Texte courant
\definecolor{GrisMoyen}{RGB}{100, 100, 100}          % Annotations
\definecolor{GrisClair}{RGB}{160, 160, 160}          % Légendes, pieds de page
\definecolor{GrisTresLeger}{RGB}{240, 240, 240}      % Fonds discrets

% ── Couleurs fonctionnelles ───────────────────────────────────────────────────
\definecolor{FondDef}{RGB}{225, 237, 252}
\definecolor{BordDef}{RGB}{0, 56, 120}
\definecolor{FondAlerte}{RGB}{255, 237, 210}
\definecolor{BordAlerte}{RGB}{210, 100, 0}
\definecolor{FondCasPratique}{RGB}{232, 244, 255}
\definecolor{BordCasPratique}{RGB}{0, 102, 179}
\definecolor{FondInfo}{RGB}{248, 248, 248}
\definecolor{BordInfo}{RGB}{130, 130, 130}
\definecolor{FondCode}{RGB}{246, 246, 246}
\definecolor{BordCode}{RGB}{0, 102, 179}

% ── Tableaux ─────────────────────────────────────────────────────────────────
\definecolor{TableauPair}{RGB}{232, 241, 252}
\definecolor{TableauImpair}{RGB}{255, 255, 255}
\definecolor{TableauEntete}{RGB}{0, 56, 120}


% ==============================================================================
% IV. TYPOGRAPHIE
% ==============================================================================

\setstretch{1.15}
\setlength{\parindent}{1.2em}
\setlength{\parskip}{0.4ex}


% ==============================================================================
% V. STYLES DE CHAPITRES ET SECTIONS
% ==============================================================================

\RequirePackage{titlesec}
\RequirePackage{titletoc}

% --- Chapitre : carré bleu marine + titre en capitales + filet orange ---
\titleformat{\chapter}[display]
    {\normalfont\bfseries\color{BleuPrimaire}}
    {%
        \filleft
        \begin{tikzpicture}
            \fill[BleuPrimaire] (0,0) rectangle (1.2,1.2);
            \node[white, font=\fontsize{28}{28}\selectfont\bfseries]
                at (0.6,0.6) {\thechapter};
        \end{tikzpicture}
    }
    {-0.5ex}
    {\filright\fontsize{18}{22}\selectfont\MakeUppercase}
    [\vspace{0.5ex}{\color{OrangeAccent}\titlerule[2pt]}\vspace{1ex}]

\titlespacing*{\chapter}{0pt}{10pt}{25pt}

% --- Section : numéro sur fond bleu marine + filet bleu secondaire ---
\titleformat{\section}
    {\normalfont\large\bfseries\color{BleuPrimaire}}
    {\colorbox{BleuPrimaire}{\color{white}\,\thesection\,}}
    {0.7em}{}
    [\vspace{-0.4ex}{\color{BleuSecondaire}\titlerule[0.8pt]}]

\titlespacing*{\section}{0pt}{1.4ex plus 0.5ex}{0.8ex plus 0.3ex}

% --- Sous-section : bleu secondaire ---
\titleformat{\subsection}
    {\normalfont\normalsize\bfseries\color{BleuSecondaire}}
    {\thesubsection}{0.6em}{}

\titlespacing*{\subsection}{0pt}{1.2ex plus 0.4ex}{0.5ex plus 0.2ex}

% --- Sous-sous-section : orange accent, italique ---
\titleformat{\subsubsection}
    {\normalfont\normalsize\itshape\color{OrangeAccent}}
    {\thesubsubsection}{0.5em}{}

\titlespacing*{\subsubsection}{0pt}{0.8ex plus 0.3ex}{0.3ex plus 0.1ex}


% ==============================================================================
% VI. EN-TÊTES ET PIEDS DE PAGE
% ==============================================================================
%
% En-tête gauche  : titre du chapitre courant (bleu)
% En-tête droite  : code UE en petites capitales (gris)  ← \docue
% Filet haut      : bleu marine
% Pied gauche     : auteur  ← \docauteur
% Pied centre     : numéro de page (bleu)
% Pied droite     : titre court  ← \doctitre
% Filet bas       : orange accent
% ==============================================================================

\pagestyle{fancy}
\fancyhf{}

\renewcommand{\headrulewidth}{0.6pt}
\renewcommand{\footrulewidth}{0.4pt}
\renewcommand{\headrule}{\color{BleuPrimaire}\hrule width\headwidth height\headrulewidth}
\renewcommand{\footrule}{\color{OrangeAccent}\hrule width\headwidth height\footrulewidth}

\fancyhead[LE,LO]{\small\color{BleuPrimaire}\leftmark}
\fancyhead[RE,RO]{\small\color{GrisMoyen}\textsc{UE\,:\,\docue}}

\fancyfoot[LE,LO]{\footnotesize\color{GrisClair}\textit{\docauteur}}
\fancyfoot[C]{\footnotesize\color{BleuPrimaire}\textbf{---\ \thepage\ ---}}
\fancyfoot[RE,RO]{\footnotesize\color{GrisClair}\textit{\doctitre}}

\fancypagestyle{plain}{%
    \fancyhf{}
    \renewcommand{\headrulewidth}{0pt}
    \renewcommand{\footrulewidth}{0.4pt}
    \renewcommand{\footrule}{\color{OrangeAccent}\hrule width\headwidth height\footrulewidth}
    \fancyfoot[C]{\footnotesize\color{BleuPrimaire}\textbf{---\ \thepage\ ---}}
    \fancyfoot[LE,LO]{\footnotesize\color{GrisClair}\textit{\docauteur}}
    \fancyfoot[RE,RO]{\footnotesize\color{GrisClair}\textit{\docinstitution}}
}


% ==============================================================================
% VII. TABLE DES MATIÈRES
% ==============================================================================

\contentsmargin{0pt}

\titlecontents{chapter}
    [1.5em]
    {\addvspace{1ex}\bfseries\color{BleuPrimaire}}
    {\contentslabel{1.5em}}
    {\hspace*{-1.5em}}
    {\hfill\contentspage}
    [\addvspace{0.3ex}]

\titlecontents{section}
    [3.8em]
    {\small\color{GrisAncre}}
    {\contentslabel{2.3em}}
    {\hspace*{-2.3em}}
    {\ \titlerule*[0.5pc]{.}\contentspage}

\titlecontents{subsection}
    [6.5em]
    {\footnotesize\color{GrisMoyen}}
    {\contentslabel{3em}}
    {\hspace*{-3em}}
    {\ \titlerule*[0.5pc]{.}\contentspage}


% ==============================================================================
% VIII. ENVIRONNEMENTS TCOLORBOX
% ==============================================================================
%
% Boîtes disponibles :
%   \begin{boxdef}      …  \end{boxdef}        Définition         (bleu pâle)
%   \begin{boxalert}    …  \end{boxalert}       Attention          (orange pâle)
%   \begin{boxinfo}     …  \end{boxinfo}        Note               (gris léger)
%   \begin{boxobjectifs}…  \end{boxobjectifs}   Objectifs          (bleu clair)
%   \begin{boxresume}   …  \end{boxresume}      À retenir          (gris + filets orange)
%   \begin{boxexo}      …  \end{boxexo}         Exercice           (blanc, filet gauche)
%   \begin{boxcas}[Titre optionnel] … \end{boxcas}  Cas pratique  (bleu très pâle)
%
% Toutes les boîtes acceptent un argument optionnel pour personnaliser les
% options tcolorbox, ex. : \begin{boxdef}[title={Ma définition}]
% ==============================================================================

% --- A. Définition ---
\newtcolorbox{boxdef}[1][]{
    enhanced, breakable,
    colback=FondDef, colframe=BordDef,
    fonttitle=\bfseries\small\color{white},
    title={D\'efinition},
    attach boxed title to top left={yshift=-2mm, xshift=5mm},
    boxed title style={colback=BleuPrimaire, colframe=BleuPrimaire, sharp corners},
    sharp corners, boxrule=0pt, leftrule=4pt,
    left=8pt, right=8pt, top=10pt, bottom=8pt,
    #1
}

% --- B. Attention / Alerte ---
\newtcolorbox{boxalert}[1][]{
    enhanced, breakable,
    colback=FondAlerte, colframe=BordAlerte,
    fonttitle=\bfseries\small\color{white},
    title={Attention},
    attach boxed title to top left={yshift=-2mm, xshift=5mm},
    boxed title style={colback=OrangeAccent, colframe=OrangeAccent, sharp corners},
    sharp corners, boxrule=0pt, leftrule=5pt,
    left=8pt, right=8pt, top=10pt, bottom=8pt,
    #1
}

% --- C. Cas pratique (titre personnalisable) ---
\newtcolorbox{boxcas}[2][]{
    enhanced, breakable,
    colback=FondCasPratique, colframe=BordCasPratique,
    fonttitle=\bfseries\small\color{white},
    title={Cas Pratique\ifx&#2&\else{} \textemdash{} #2\fi},
    attach boxed title to top left={yshift=-2mm, xshift=5mm},
    boxed title style={colback=BleuSecondaire, colframe=BleuSecondaire, sharp corners},
    sharp corners, boxrule=0.8pt,
    left=8pt, right=8pt, top=10pt, bottom=8pt,
    #1
}

% --- D. Note / Information ---
\newtcolorbox{boxinfo}[1][]{
    enhanced, breakable,
    colback=FondInfo, colframe=BordInfo,
    fonttitle=\bfseries\small\color{GrisAncre},
    title={Note},
    attach boxed title to top left={yshift=-2mm, xshift=5mm},
    boxed title style={colback=GrisTresLeger, colframe=BordInfo, sharp corners},
    sharp corners, boxrule=0pt, leftrule=2pt,
    left=8pt, right=8pt, top=10pt, bottom=8pt,
    #1
}

% --- E. Objectifs de chapitre ---
\newtcolorbox{boxobjectifs}[1][]{
    enhanced, breakable,
    colback=BleuClair, colframe=BleuPrimaire,
    fonttitle=\bfseries\small\color{white},
    title={Objectifs du chapitre},
    attach boxed title to top left={yshift=-2mm, xshift=5mm},
    boxed title style={colback=BleuPrimaire, colframe=BleuPrimaire, sharp corners},
    sharp corners, boxrule=0.8pt,
    left=8pt, right=8pt, top=10pt, bottom=8pt,
    #1
}

% --- F. À retenir ---
\newtcolorbox{boxresume}[1][]{
    enhanced, breakable,
    colback=GrisTresLeger, colframe=OrangeAccent,
    fonttitle=\bfseries\small\color{white},
    title={\`A retenir},
    attach boxed title to top left={yshift=-2mm, xshift=5mm},
    boxed title style={colback=OrangeAccent, colframe=OrangeAccent, sharp corners},
    sharp corners, boxrule=0pt, toprule=2pt, bottomrule=2pt,
    left=8pt, right=8pt, top=10pt, bottom=8pt,
    #1
}

% --- G. Exercice ---
\newtcolorbox{boxexo}[1][]{
    enhanced, breakable,
    colback=white, colframe=BleuSecondaire,
    fonttitle=\bfseries\small\color{white},
    title={Exercice},
    attach boxed title to top left={yshift=-2mm, xshift=5mm},
    boxed title style={colback=BleuSecondaire, colframe=BleuSecondaire, sharp corners},
    sharp corners, boxrule=0pt, leftrule=3pt,
    left=8pt, right=8pt, top=10pt, bottom=8pt,
    #1
}


% ==============================================================================
% IX. STYLE DE CODE SOURCE
% ==============================================================================

\lstdefinestyle{StyleCode}{
    backgroundcolor=\color{FondCode},
    basicstyle=\ttfamily\footnotesize\color{GrisAncre},
    keywordstyle=\bfseries\color{BleuPrimaire},
    commentstyle=\itshape\color{GrisClair},
    stringstyle=\color{OrangeAccent},
    numberstyle=\tiny\color{GrisClair},
    breakatwhitespace=false, breaklines=true,
    captionpos=b, keepspaces=true,
    numbers=left, numbersep=8pt,
    showspaces=false, showstringspaces=false, showtabs=false,
    tabsize=4,
    frame=single, frameround=ffff,
    rulecolor=\color{BordCode},
    autogobble=true,
    xleftmargin=12pt, xrightmargin=4pt,
    aboveskip=8pt, belowskip=6pt,
    literate=
        {é}{{\'e}}1 {è}{{\`e}}1 {ê}{{\^e}}1 {ë}{{\"e}}1
        {à}{{\`a}}1 {â}{{\^a}}1 {ç}{{\c{c}}}1 {ù}{{\`u}}1
        {û}{{\^u}}1 {î}{{\^i}}1 {ô}{{\^o}}1
}

\lstset{style=StyleCode}

\lstdefinelanguage{SQL}{
    keywords={SELECT, FROM, WHERE, AND, OR, NOT, INSERT, INTO, VALUES,
              UPDATE, SET, DELETE, CREATE, TABLE, ALTER, DROP, INDEX,
              JOIN, LEFT, RIGHT, INNER, OUTER, ON, GROUP, BY, ORDER,
              HAVING, AS, DISTINCT, UNION, ALL, IN, EXISTS, LIKE,
              IS, NULL, BETWEEN, CASE, WHEN, THEN, ELSE, END,
              PRIMARY, KEY, FOREIGN, REFERENCES, UNIQUE, CHECK,
              DATABASE, USE, SHOW, DESCRIBE, GRANT, REVOKE},
    sensitive=false,
    morecomment=[l]{--},
    morecomment=[s]{/*}{*/},
    morestring=[b]', morestring=[b]"
}


% ==============================================================================
% X. LISTES PERSONNALISÉES
% ==============================================================================

\setlist[itemize,1]{
    leftmargin=1.8em,
    label={\color{BleuSecondaire}\textbullet},
    itemsep=0.3ex
}
\setlist[itemize,2]{
    leftmargin=1.5em,
    label={\color{OrangeAccent}$\triangleright$},
    itemsep=0.2ex
}
\setlist[itemize,3]{
    leftmargin=1.2em,
    label={\color{GrisMoyen}$\circ$},
    itemsep=0.1ex
}
\setlist[enumerate,1]{
    leftmargin=2em,
    label={\color{BleuPrimaire}\bfseries\arabic*.},
    itemsep=0.3ex
}
\setlist[enumerate,2]{
    leftmargin=1.8em,
    label={\color{BleuSecondaire}\alph*)},
    itemsep=0.2ex
}


% ==============================================================================
% XI. ENVIRONNEMENT TABLEAU ACADÉMIQUE
% ==============================================================================
%
% Usage :
%   \begin{tableaucours}{<spec colonnes>}{<en-tête>}
%     contenu …
%   \end{tableaucours}
%
% Exemple :
%   \begin{tableaucours}{>{\bfseries}p{3cm} p{10cm}}{Terme & Définition}
%     Risque & Probabilité d'occurrence d'une menace … \\
%   \end{tableaucours}
% ==============================================================================

\newenvironment{tableaucours}[2]{%
    \begin{table}[H]
    \centering\small
    \rowcolors{2}{TableauPair}{TableauImpair}
    \begin{tabular}{#1}
    \toprule
    \rowcolor{TableauEntete}
    \color{white}\bfseries #2
    \midrule
}{%
    \bottomrule
    \end{tabular}
    \end{table}
}


% ==============================================================================
% XII. COMMANDES UTILITAIRES
% ==============================================================================

% Terme clé indexé et mis en gras bleu
\newcommand{\termecle}[1]{%
    \textbf{\color{BleuPrimaire}#1}\index{#1}%
}

% Sigle en petites capitales
\newcommand{\sigle}[1]{\textsc{#1}}

% Note en marge
\newcommand{\notemarge}[1]{%
    \marginpar{\footnotesize\color{GrisMoyen}\itshape #1}%
}

% Niveaux de difficulté
\newcommand{\niveaubase}{%
    \textcolor{GrisMoyen}{\small\textsc{[Niveau\,1]}}}
\newcommand{\niveaumoyen}{%
    \textcolor{BleuSecondaire}{\small\textsc{[Niveau\,2]}}}
\newcommand{\niveauexpert}{%
    \textcolor{OrangeAccent}{\small\textsc{[Niveau\,3]}}}

% Référence à une norme (ex. : \norme{ISO 27001})
\newcommand{\norme}[1]{\textbf{\color{OrangeAccent}[#1]}}

% Séparateur décoratif entre sections
\newcommand{\separateur}{%
    \vspace{1ex}
    \begin{center}
        {\color{BleuMoyen}\rule{0.3\linewidth}{0.5pt}}%
        \quad{\color{OrangeAccent}\rule{0.06\linewidth}{1pt}}\quad%
        {\color{BleuMoyen}\rule{0.3\linewidth}{0.5pt}}
    \end{center}
    \vspace{0.6ex}
}


% ==============================================================================
% XIII. PAGE DE TITRE — \maketitlecours
% ==============================================================================
%
% Place \maketitlecours juste après \begin{document}.
% Elle utilise automatiquement les métadonnées définies en section II.
%
% Structure :
%   — Double cadre (noir fin extérieur + orange épais intérieur)
%   — Tableau bilingue FR | Logo | EN  ← logo fourni via \doclogo
%   — Bandeau de type de document : \doctypedevoir — UE : \docue
%   — Titre principal : \doctitre
%   — Sous-titre : \docsoustitre
%   — Bloc normes (optionnel) : \docnormes  (vide = non affiché)
%   — Épigraphe (optionnel) : \docepigraphe (vide = non affiché)
%   — Bloc auteur : \docauteur / \docinstitution / \docdepartement / \docanneescolaire
% ==============================================================================

\newcommand{\maketitlecours}{%
    \thispagestyle{empty}
    \begin{titlepage}

    % ── Double cadre ─────────────────────────────────────────────
    \begin{tikzpicture}[remember picture, overlay]
        \draw[NoirProfond, line width=1pt]
            ([xshift=0.7cm,  yshift=-0.7cm]  current page.north west)
            rectangle
            ([xshift=-0.7cm, yshift=0.7cm]   current page.south east);
        \draw[OrangeAccent, line width=2.5pt]
            ([xshift=0.95cm,  yshift=-0.95cm]  current page.north west)
            rectangle
            ([xshift=-0.95cm, yshift=0.95cm]   current page.south east);
    \end{tikzpicture}

    \vspace*{0.3cm}
    \begin{adjustwidth}{0.6cm}{0.6cm}

    % ── Tableau bilingue FR | Logo | EN ──────────────────────────
    {\renewcommand{\arraystretch}{0.45}
    \noindent
    \begin{tabular}{@{}
        >{\centering\arraybackslash\footnotesize}p{0.34\linewidth}
        >{\centering\arraybackslash}p{0.22\linewidth}
        >{\centering\arraybackslash\footnotesize}p{0.34\linewidth}
        @{}}
        \textbf{REPUBLIQUE DU CAMEROUN}
            & \multirow{16}{*}{\includegraphics[width=3.0cm]{\doclogo}}
            & \textbf{REPUBLIC OF CAMEROON}          \\
        Paix-Travail-Patrie          & & Peace - Work - Fatherland\\
        *****     & & *****\\
        \textbf{MINIST\`ERE DE L'ENSEIGNEMENT SUP\'ERIEUR}
            & & \textbf{MINISTRY OF HIGHER EDUCATION} \\
        ***** & & *****                               \\
        \textbf{UNIVERSIT\'E DE YAOUND\'E I}
            & & \textbf{UNIVERSITY OF YAOUND\'E I}    \\
        ***** & & *****                               \\
        \textbf{\'ECOLE NATIONALE SUP\'ERIEURE POLYTECHNIQUE}
            & & \textbf{NATIONAL ADVANCED SCHOOL OF ENGINEERING} \\
        ***** & & *****                               \\
        \textbf{D\'EPARTEMENT DE G\'ENIE INFORMATIQUE}
            & & \textbf{COMPUTER ENGINEERING DEPARTMENT} \\
        ***** & & *****
    \end{tabular}}

    \vspace{0.8cm}

    % ── Type de document ─────────────────────────────────────────
    \begin{center}
        {\color{GrisMoyen}\small\textsc{%
            \doctypedevoir{} --- UE\,:\,\docue%
        }}\\[1.2ex]
        {\color{BleuMoyen}\rule{0.5\linewidth}{0.5pt}}
    \end{center}

    \vspace{0.5cm}

    % ── Titre principal ──────────────────────────────────────────
    \begin{center}
        {\color{BleuPrimaire}
            \fontsize{22}{28}\selectfont\bfseries
            \MakeUppercase{\doctitre}
        }\\[1.0ex]
        {\color{OrangeAccent}\large\itshape \docsoustitre}
    \end{center}

    \vspace{0.6cm}

    % ── Filet médian ─────────────────────────────────────────────
    \begin{center}
        {\color{BleuMoyen}\rule{0.28\linewidth}{0.5pt}}%
        \enspace{\color{OrangeAccent}\rule{0.08\linewidth}{1.5pt}}\enspace%
        {\color{BleuMoyen}\rule{0.28\linewidth}{0.5pt}}
    \end{center}

    \vspace{0.5cm}

    % ── Normes (optionnel — vide si \docnormes est vide) ─────────
    \ifx\docnormes\empty\else
    \begin{center}
        {\color{GrisClair}\footnotesize \docnormes}
    \end{center}
    \vspace{0.4cm}
    \fi

    % ── Épigraphe (optionnel — vide si \docepigraphe est vide) ───
    \ifx\docepigraphe\empty\else
    \vspace{0.3cm}
    \begin{center}
        {\color{GrisMoyen}\itshape\small \og{}\docepigraphe\fg{}}
    \end{center}
    \vspace{0.3cm}
    \fi

    \vfill

    % ── Bloc auteur ──────────────────────────────────────────────
    {\color{white}
        \colorbox{OrangeAccent}{%
            \begin{minipage}{\linewidth}
                \vspace{0.5cm}
                \begin{center}
                    {\large\bfseries\docauteur}\\[0.5ex]
                    {\small\itshape \docinstitution\ ---\ \docdepartement}\\[0.6ex]
                    {\normalsize Ann\'ee acad\'emique \docanneescolaire}
                \end{center}
                \vspace{0.5cm}
            \end{minipage}%
        }%
    }

    \end{adjustwidth}
    \end{titlepage}
    \setcounter{page}{1}
}


% ==============================================================================
% XIV. ENVIRONNEMENTS MATHÉMATIQUES
% ==============================================================================

\theoremstyle{definition}
\newtheorem{theoreme}{Th\'eor\`eme}[chapter]
\newtheorem{definition}{D\'efinition}[chapter]
\newtheorem{proposition}{Proposition}[chapter]
\newtheorem{corollaire}{Corollaire}[chapter]
\newtheorem{lemme}{Lemme}[chapter]

\theoremstyle{remark}
\newtheorem{remarque}{Remarque}[chapter]
\newtheorem{exemple}{Exemple}[chapter]


% ==============================================================================
% XV. PARAMÈTRES FINAUX
% ==============================================================================

\setcounter{secnumdepth}{3}
\setcounter{tocdepth}{3}
\hyphenpenalty=800
\exhyphenpenalty=800
\color{GrisAncre}

\endinput
% ==============================================================================
% Fin de style.sty — MAGUENA NDENE Marcel Allan
% ==============================================================================